\section{预训练目标构造}
\label{sec:nlp-llm-objectives}

前几节建立了从词元序列到词表概率的计算，但这个概率函数本身没有规定应当学会什么。同一份文本可以要求模型只根据左侧续写，也可以要求它根据缺口两侧恢复中间片段；同一个网络既能对每个位置计分，也能只对少数被遮蔽的位置计分。构造预训练目标，就是为每个训练样本规定三件事：预测时允许看到哪些位置，要预测哪些位置，以及每个预测的错误按多大权重计入损失。

这一构造面临三个相互牵连的挑战。其一，条件和预测对象一旦改变，模型拟合的就是另一个条件分布，学到的能力也随之不同。其二，目标必须与注意力掩码、网络结构和生成方式配套，否则要么在训练时泄漏答案，要么训练出的模型无法直接用于从左到右生成。其三，不同目标在训练信号的密度、能否利用双侧上下文和能否直接生成之间各有取舍。

围绕这些挑战，常见目标可以归为三类。自回归类只让模型看左侧，包括下一词元预测、在不改结构的前提下获得填空能力的中间填充，以及为每个位置增加远期监督的多词元预测。去噪重建类先损坏输入再要求恢复，可见条件包含两侧上下文，包括遮蔽语言模型、T5的跨度损坏、BART的去噪、前缀语言模型以及把它们混合起来的UL2。掩码扩散类把去噪推广为从全掩码序列逐步生成的过程。最后，目标与结构相互约束，我们比较与这些目标配套的编码器—解码器结构和仅解码器结构，并在节末用同一条通知把各目标的输入、掩码与目标序列并列写出。

\subsection{可见条件与预测对象的统一描述}
\label{sec:nlp-llm-objective-visibility}
\label{sec:nlp-llm-objective-loss}

先固定一条贯穿本节的教学序列：通知“读书会／改到／周五／下午／两点／举行”经教学词元化后得到六个词元$x_{1:6}$。若要学习从左到右续写，预测$x_3=$“周五”时只能看到“读书会改到”；若要学习填空，可以遮住“周五”而同时提供两侧的“读书会改到”和“下午两点举行”。两种做法预测的是同一个词元，但待学的条件分布不同：
$p_\theta(x_3\mid x_1,x_2)$要求模型在不知道后文时猜测改期日期，可能的取值很多；$p_\theta(x_3\mid x_1,x_2,x_4,x_5,x_6)$则已经知道是“下午两点”，候选少得多。条件更多时预测更容易，损失也更低，但这个较低的损失回答的是另一个问题，不能代替前者。

为了把不同目标放在同一框架下比较，我们把一个训练样本的目标描述为三元组。令$\mathcal S(x)$为参与计分的\emph{预测位置集合}；对每个$i\in\mathcal S(x)$，令$V_i(x)$为预测$x_i$时允许读取的\emph{可见位置集合}；令$w_i(x)\geq0$为该预测的计分权重。样本连同目标构造共同产生一组训练记录$\{(x_{V_i(x)},x_i,w_i(x)):i\in\mathcal S(x)\}$。下一词元预测取$\mathcal S=\{1,\ldots,T\}$、$V_i=\{1,\ldots,i-1\}$；遮蔽类目标中$\mathcal S$只含被遮蔽位置，可见集合为全部位置，但被遮蔽位置的内容已被替换为掩码符号。开始、结束和填充等特殊符号是否属于$\mathcal S$，也在这一步规定。

在Transformer中，可见集合通过\term{注意力掩码}（attention mask）实现。注意力掩码是一个二值矩阵$\bm M\in\{0,1\}^{T\times T}$，$M_{i,j}=1$表示位置$i$的查询可以读取位置$j$的键和值。计算注意力分数时，对$M_{i,j}=0$的位置加上$-\infty$，softmax后其权重为零，归一化只在允许位置上进行。图~\ref{fig:nlp-llm-objective-masks}把本节将要用到的四种掩码画在同一个六位置网格上：因果掩码是下三角，前缀语言模型在前缀块内放开，双向掩码全部放开，编码器—解码器则把源和目标分成两块并经交叉注意力相连。注意力掩码只限制信息流，不标记哪些位置计分；计分由$\mathcal S$和$w_i$决定，两者需要分别设置。

\begin{figure*}[htbp]
  \centering
  % 第九卷 3.6 节：四种目标配套的可见范围（注意力掩码）对照。
% 行为查询位置 i，列为键位置 j；六个位置依次为“读书会／改到／周五／下午／两点／举行”。
% 蓝格：同一组参数内的自注意力可见；黄格：编码器—解码器中经交叉注意力读取；白格：不可见。
\begin{tikzpicture}[
  x=1cm,
  y=1cm,
  font=\figurefont\footnotesize,
  seen/.style={draw=white,line width=0.5pt,fill=FigureBlue!70},
  cross/.style={draw=white,line width=0.5pt,fill=FigureYellow!85},
  hidden/.style={draw=FigureGrid,line width=0.4pt,fill=white},
  idx/.style={font=\figurefont\scriptsize,text=FigureMuted},
  ptitle/.style={font=\figurefont\bfseries\footnotesize,text=FigureInk,anchor=south},
  note/.style={font=\figurefont\scriptsize,text=FigureMuted,anchor=north,align=center}
]
  \def\s{0.40}
  % 面板编号 \p：1 因果，2 前缀（P=2），3 双向，4 编码器—解码器（源为位置 1–3）。
  \foreach \p/\xo/\ttl/\nt in {%
      1/0.0/{(a) 因果（下一词元预测）}/{位置$i$只读$j\leq i$},%
      2/3.55/{(b) 前缀语言模型}/{前缀$P=2$内互读，\\其后因果},%
      3/7.1/{(c) 双向（掩码语言模型）}/{全部互读；\\被遮位置输入为掩码},%
      4/10.65/{(d) 编码器—解码器}/{源1--3双向；目标4--6因果，\\并经交叉注意力读源}} {
    \foreach \i in {1,...,6} {
      \foreach \j in {1,...,6} {
        \pgfmathtruncatemacro{\kind}{%
          \p==1 ? (\j<=\i ? 1 : 0) :
          (\p==2 ? ((\j<=2 || \j<=\i) ? 1 : 0) :
          (\p==3 ? 1 :
          (\i<=3 ? (\j<=3 ? 1 : 0) : (\j<=3 ? 2 : (\j<=\i ? 1 : 0)))))}
        \ifcase\kind
          \path[hidden] ({\xo+(\j-1)*\s},{-(\i-1)*\s}) rectangle ++(\s,-\s);
        \or
          \path[seen] ({\xo+(\j-1)*\s},{-(\i-1)*\s}) rectangle ++(\s,-\s);
        \or
          \path[cross] ({\xo+(\j-1)*\s},{-(\i-1)*\s}) rectangle ++(\s,-\s);
        \fi
      }
      \node[idx,anchor=east] at ({\xo-0.04},{-(\i-0.5)*\s}) {\i};
      \node[idx,anchor=south] at ({\xo+(\i-0.5)*\s},0.02) {\i};
    }
    \draw[FigureStroke,line width=0.6pt] (\xo,0) rectangle ++({6*\s},{-6*\s});
    \node[ptitle] at ({\xo+3*\s},0.36) {\ttl};
    \node[note] at ({\xo+3*\s},{-6*\s-0.12}) {\nt};
  }
  % 前缀与源的分界线。
  \draw[FigureStroke,line width=0.9pt,dashed] ({3.55+2*\s},0) -- ({3.55+2*\s},{-6*\s});
  \draw[FigureStroke,line width=0.9pt,dashed] (3.55,{-2*\s}) -- ({3.55+6*\s},{-2*\s});
  \draw[FigureStroke,line width=0.9pt] ({10.65+3*\s},0) -- ({10.65+3*\s},{-6*\s});
  \draw[FigureStroke,line width=0.9pt] (10.65,{-3*\s}) -- ({10.65+6*\s},{-3*\s});
  % 坐标说明与图例。
  \node[idx,anchor=east] at (-0.3,{-3*\s}) {查询$i$};
  \node[idx,anchor=south east] at (-0.06,0.02) {键$j$};
  \path[seen] (3.2,-3.75) rectangle ++(\s,-\s);
  \node[anchor=west,font=\figurefont\scriptsize,text=FigureInk] at ({3.2+\s+0.08},{-3.75-0.5*\s}) {自注意力可见};
  \path[cross] (6.0,-3.75) rectangle ++(\s,-\s);
  \node[anchor=west,font=\figurefont\scriptsize,text=FigureInk] at ({6.0+\s+0.08},{-3.75-0.5*\s}) {交叉注意力读取};
  \path[hidden] (8.9,-3.75) rectangle ++(\s,-\s);
  \node[anchor=west,font=\figurefont\scriptsize,text=FigureInk] at ({8.9+\s+0.08},{-3.75-0.5*\s}) {不可见};
\end{tikzpicture}

  \caption{四种预训练目标配套的可见范围。每个网格的行是查询位置$i$，列是键位置$j$，六个位置依次对应教学序列“读书会／改到／周五／下午／两点／举行”；有色格表示位置$i$可以读取位置$j$。(a)因果掩码为下三角；(b)前缀语言模型取前缀长度$P=2$，虚线左上块内双向，其余行仍为因果；(c)双向掩码全部可见，被遮蔽位置以掩码符号作为输入；(d)编码器—解码器以位置1--3为源、4--6为目标，黄色块经由单独参数的交叉注意力读取编码器最后一层输出。(b)与(d)的可见模式相同，区别在于参数是否分开以及目标侧读取源表示的方式。}
  \label{fig:nlp-llm-objective-masks}
\end{figure*}

有了三元组，各种目标都可以写成同一形式的加权负对数似然。若一次预测给真实目标分配的概率为$p$，损失$-\log p$在$p$接近1时趋于零，在$p$接近0时迅速增大。对一个样本及其随机损坏，一般的计分形式为
\begin{equation}
 \ell(\theta;x)=
 -\frac{1}{Z(x)}\sum_{i\in\mathcal S(x)}
 w_i(x)\log p_\theta\bigl(x_i\mid x_{V_i(x)}\bigr),
 \label{eq:nlp-llm-objective-general}
\end{equation}
其中$Z(x)>0$是归一化常数。式\eqref{eq:nlp-llm-objective-general}把一个目标拆成三个可以独立检查的选择：$V_i$规定条件，$\mathcal S$规定预测对象，$w_i$和$Z$规定各项在总目标中的分量。后文介绍的每种目标都只是这三项的不同取法，配上相应的样本构造。

归一化常数的取法会改变长短样本的相对权重。若$w_i=1$且$Z(x)=|\mathcal S(x)|$，每个样本先求平均再在批次内平均；若把一批样本的所有计分词元汇总后再平均，每个词元等权，长样本的贡献随之增大。以教学数值为例：样本甲有2个计分词元、平均损失1.0，样本乙有8个计分词元、平均损失2.0。按样本平均得$(1.0+2.0)/2=1.5$，按词元平均得$(2\times1.0+8\times2.0)/10=1.8$。两种口径都常见，报告平均损失时需要写明采用哪一种。

对数据分布与目标构造中的随机性取期望，就得到总体目标$\mathcal L(\theta)=\mathbb E[\ell(\theta;x)]$，训练用小批次估计它的梯度。这个目标拟合的是什么？固定一种条件$C$，设数据中目标词元的条件分布为$q(\cdot\mid C)$，则等权的期望交叉熵可以分解为
\[
 -\mathbb E_{y\sim q(\cdot\mid C)}\log p_\theta(y\mid C)
 =H\bigl(q(\cdot\mid C)\bigr)+D_{\mathrm{KL}}\bigl(q(\cdot\mid C)\,\|\,p_\theta(\cdot\mid C)\bigr).
\]
第一项与参数无关，因此最小化交叉熵等价于让模型的条件分布逼近数据的条件分布。这说明目标的选择决定了“逼近哪一个条件分布”：在模型表达能力充分、能求得最优值的理想情形下，下一词元预测逼近$q(x_i\mid x_{<i})$，填空目标逼近$q(x_i\mid\text{两侧})$。实际训练只有有限数据与有限优化，这个解释给出的是目标指向的方向，不是对任意新输入的保证。

泄漏不只发生在最后一层。多层注意力中，位置$i$第$l$层的表示依赖于它在第$l-1$层可读位置的表示，而后者又依赖更早一层可读的位置。因此位置$i$最终能够获取的信息，是掩码图中从$i$出发经不超过$L$步可以到达的全部位置。若某个中间层让位置2读取了位置3，那么后续任何读取位置2的查询，都可能间接得到$x_3$。因果掩码和前缀掩码的可达关系是传递闭合的，多层叠加后可见范围不会扩大；自行设计掩码时，则需要对整个计算图检查这种可达性。合法的复制不属于泄漏，例如从可见材料中抄出实体名称；需要排除的是目标经由本应不可见的路径进入预测。

由式\eqref{eq:nlp-llm-objective-general}还可以看出，不同目标得到的损失值一般不能直接比较。条件集合、计分位置、归一化口径和词元化中任意一项不同，同一个数值就对应不同的问题。比较目标时，应在固定的评测协议下比较下游行为，而不是比较各自的训练损失。

\subsection{基于自回归预测的目标}
\label{sec:nlp-llm-objective-autoregressive}

自回归目标让每个位置只看左侧。它与逐词元生成的过程一致：训练时学的条件分布，正是生成时每一步要用的条件分布。由于所有位置都可以计分，一次前向计算就为长度为$T$的样本提供$T$个监督信号。这类目标的局限有两点：它不能利用右侧上下文完成填空，而代码补全等应用恰好需要光标两侧的内容；它只为下一个位置提供监督，模型可能更多依赖局部搭配，而较少为更远的未来做准备。中间填充通过重排文档解决前一个问题，而不改变网络结构；多词元预测通过增加输出模块缓解后一个问题。三者共享同一个因果解码器，差别只在样本构造与输出头。

\subsubsection{下一词元预测}
\label{sec:nlp-llm-objective-positions}

\term{下一词元预测}（next-token prediction, NTP）对每个位置预测其后继，学习目标为
\begin{equation}
 \mathcal L_{\mathrm{NTP}}(\theta)=
 -\sum_{t=1}^{T}\log p_\theta(x_t\mid x_{<t}),
 \label{eq:nlp-llm-ntp}
\end{equation}
它正是式\eqref{eq:nlp-llm-sequence-probability}中序列概率的负对数，因此最小化它就是对训练文本做最大似然估计。在三元组中，$\mathcal S=\{1,\ldots,T\}$，$V_t=\{1,\ldots,t-1\}$，$w_t=1$。

实现时，输入序列与目标序列错开一位。把开始标记记为BOS，网络读入$(\mathrm{BOS},x_1,\ldots,x_{T-1})$，第$t$个输出位置的分布与目标$x_t$比较。因果掩码保证第$t$个输出只读取BOS与$x_{1:t-1}$。训练时条件中的前缀总是真实文本，而不是模型自己生成的词元，这种做法即第\ref{sec:nlp-llm-probabilities}小节的教师强制；它让$T$个位置的预测在一次前向计算中并行完成，代价是训练时的前缀分布与生成时不同，相关讨论见第\ref{chap:nlp-supervised-finetuning}章。对教学序列，下表列出每条记录的条件与目标：
\begin{center}
 \begin{tabular}{@{}cll@{}}
  \toprule
  位置$t$ & 条件$x_{<t}$ & 目标$x_t$\\
  \midrule
  1 & BOS & 读书会\\
  2 & BOS 读书会 & 改到\\
  3 & BOS 读书会 改到 & 周五\\
  4 & BOS 读书会 改到 周五 & 下午\\
  5 & BOS 读书会 改到 周五 下午 & 两点\\
  6 & BOS 读书会 改到 周五 下午 两点 & 举行\\
  \bottomrule
 \end{tabular}
\end{center}
预测“周五”时看不到“下午两点”，预测“举行”时可以看到全部前文。一个六词元样本产生了长度从1到6的六种前缀，这正是自回归目标训练信号密集的来源。

实际训练常把多篇文档拼接后切成定长样本（见第\ref{sec:nlp-llm-data-packing}小节），于是需要单独规定哪些位置只作条件、不计分。为此在式\eqref{eq:nlp-llm-ntp}的每一项乘以计分掩码$m_t\in\{0,1\}$，并以$Z=\sum_t m_t$归一化。填充位置取$m_t=0$；若采用文档内注意力掩码，第二篇文档的第一个词元只以本文档的开始标记为条件，不能读取前一篇的尾部；对重叠切分的长文档，重叠部分只提供条件，取$m_t=0$，只对新引入的词元计分，避免同一段文字被重复计分。改变计分位置会改变每段材料的实际学习权重，即使输入字符串完全相同。

下一词元预测的优点是信号密集、与生成方式一致、实现简单，并且每个位置的损失有明确的概率意义，可以直接计算困惑度。它的不足也来自同一个约束：条件只有左侧，既不能直接完成需要右侧信息的填空，也只对一步之后的词元负责。

\subsubsection{中间填充}
\label{sec:nlp-llm-objective-fim}

代码编辑器中的补全请求通常发生在文件中间：光标之前有函数签名，光标之后还有返回语句和调用代码。只学过下一词元预测的模型只能利用光标之前的内容。改用双向结构可以读取两侧，却会失去从左到右生成的能力。\term{中间填充}（fill-in-the-middle, FIM）的思路是不改结构、只改样本：把文档中被挖去的中间片段移到序列末尾，因果模型生成它时，两侧内容都已经出现在左边。

具体构造如下\scite{nlp-llm-fim2022}。把一篇文档在两个随机位置切开，得到前缀$\bm p$、中间$\bm m$和后缀$\bm s$，原文为$\bm p\bm m\bm s$。引入三个特殊词元$\langle\mathrm{PRE}\rangle$、$\langle\mathrm{SUF}\rangle$、$\langle\mathrm{MID}\rangle$标记各段，以及结束词元$\langle\mathrm{EOT}\rangle$。按前缀—后缀—中间（prefix-suffix-middle, PSM）顺序，训练序列为
\[
 \langle\mathrm{PRE}\rangle\,\bm p\,\langle\mathrm{SUF}\rangle\,\bm s\,\langle\mathrm{MID}\rangle\,\bm m\,\langle\mathrm{EOT}\rangle .
\]
随后照常用式\eqref{eq:nlp-llm-ntp}训练。生成$\bm m$的每个词元时，因果掩码允许读取它之前的$\bm p$和$\bm s$，于是得到了双侧条件。另一种后缀—前缀—中间（suffix-prefix-middle, SPM）顺序写为$\langle\mathrm{PRE}\rangle\langle\mathrm{SUF}\rangle\,\bm s\,\langle\mathrm{MID}\rangle\,\bm p\,\bm m\,\langle\mathrm{EOT}\rangle$，此时$\bm p\bm m$按原文顺序相连，中间片段紧接前缀生成，形式上更接近普通续写；原论文建议两种顺序联合训练。推断时，把光标前后的内容按同样格式放入$\langle\mathrm{PRE}\rangle\bm p\langle\mathrm{SUF}\rangle\bm s\langle\mathrm{MID}\rangle$，让模型生成直到$\langle\mathrm{EOT}\rangle$。

以教学序列为例，取$\bm p=$“读书会改到”、$\bm m=$“周五”、$\bm s=$“下午两点举行”，PSM序列为“$\langle\mathrm{PRE}\rangle$读书会改到$\langle\mathrm{SUF}\rangle$下午两点举行$\langle\mathrm{MID}\rangle$周五$\langle\mathrm{EOT}\rangle$”。预测“周五”时的条件已经包含“下午两点举行”，这对应前文的填空条件分布，而网络和掩码与下一词元预测完全相同。

\Needspace{12\baselineskip}
\begin{algorithm}[中间填充的样本构造]
  \label{alg:nlp-llm-fim}
  \begin{algorithmic}[1]
    \Require 文档字符串$d$，变换比例$\rho\in[0,1]$，SPM所占比例$\rho_{\mathrm{spm}}$
    \Ensure 一个用于下一词元预测的训练序列
    \State 以概率$1-\rho$返回原文$\mathrm{tok}(d)$，结束
    \State 在字符位置上均匀抽取两个切点$a\leq b$，令$\bm p=d_{<a}$，$\bm m=d_{a:b}$，$\bm s=d_{\geq b}$
    \State 分别词元化$\bm p,\bm m,\bm s$
    \If{以概率$\rho_{\mathrm{spm}}$选中SPM}
      \State 返回$\langle\mathrm{PRE}\rangle\langle\mathrm{SUF}\rangle\,\mathrm{tok}(\bm s)\,\langle\mathrm{MID}\rangle\,\mathrm{tok}(\bm p)\,\mathrm{tok}(\bm m)\,\langle\mathrm{EOT}\rangle$
    \Else
      \State 返回$\langle\mathrm{PRE}\rangle\,\mathrm{tok}(\bm p)\,\langle\mathrm{SUF}\rangle\,\mathrm{tok}(\bm s)\,\langle\mathrm{MID}\rangle\,\mathrm{tok}(\bm m)\,\langle\mathrm{EOT}\rangle$
    \EndIf
  \end{algorithmic}
\end{algorithm}

算法~\ref{alg:nlp-llm-fim}在字符层面切分，再分别词元化。若直接在词元边界切分，推断时用户光标落在一个词元中间的情形在训练中从未出现，模型会对不完整的词元前缀处理失当。变换可以在拼接切块之前对整篇文档施加，也可以在切出的训练上下文内施加；后者保证三段都落在同一个样本中，不会因切块丢掉后缀。比例$\rho$决定有多少样本被变换。Bavarian等的实验观察到，在$\rho$取到0.5甚至更高时，模型从左到右续写的测试损失几乎不受影响，同时获得了填充能力；他们把这一现象称为“填充能力几乎免费”，这一结论来自其代码与文本数据上的实验\scite{nlp-llm-fim2022}。

从三元组看，FIM没有改变$V_t$的形式，它改变的是原文位置与计算顺序之间的对应：计算序列中的“左侧”包含了原文的右侧。它可以看作只损坏一个跨度、并用因果解码器恢复的去噪目标，与后文T5的跨度损坏相比，它用重排代替编码器，因此保留了从左到右生成的能力。它的代价是需要额外的特殊词元和格式约定，且每个样本只挖去一段；被挖片段之后的$\bm m$部分预测较难，而前缀$\bm p$与后缀$\bm s$部分仍是普通续写。

\subsubsection{多词元预测}
\label{sec:nlp-llm-objective-mtp}

下一词元预测在每个位置只为一步之后负责。模型因此可能主要学习“看到前一两个词元就猜下一个”的局部规律，而对需要提前规划的内容（例如代码中稍后才闭合的结构）缺少直接的训练信号。\term{多词元预测}（multi-token prediction, MTP）让同一位置的表示同时承担更远未来的预测，以增加训练信号并促使表示编码更长程的信息。

最直接的做法是在共享主干上并列多个独立输出头\scite{nlp-llm-mtp2024}。令共享主干在位置$t$给出表示$\bm h_t=f_{\mathrm s}(x_{1:t})\in\mathbb R^{d}$，第$k$个头$f_k$（$k=1,\ldots,n$，各为一个Transformer层）把它变换为$f_k(\bm h_t)$，再经共享的输出矩阵$\bm U\in\mathbb R^{V\times d}$得到$p_{\theta,k}(x_{t+k}\mid x_{1:t})=\operatorname{softmax}(\bm U f_k(\bm h_t))$。学习目标为各头损失的加权和：
\begin{equation}
 \mathcal L_{\mathrm{MTP}}=
 -\sum_{k=1}^{n}\lambda_k
 \sum_{t:\,t+k\leq T}
 \log p_{\theta,k}(x_{t+k}\mid x_{1:t}),
 \label{eq:nlp-llm-mtp}
\end{equation}
其中$\lambda_k\geq0$为权重，$k=1$的头就是普通的下一词元预测。所有头读取的都是同一个$\bm h_t$，因此第$k$步的误差会回传到主干，要求$\bm h_t$同时携带预测$x_{t+1},\ldots,x_{t+n}$所需的信息。为控制显存，实现时可以依次计算各头的前向与反向、把梯度累加到主干，而不必同时保存$n$个大小为$V$的对数几率向量。Gloeckle等报告，这一收益随模型规模增大而更明显，在代码任务上更突出：其13B参数模型以$n=4$训练时，在HumanEval上比同等规模的下一词元预测模型多解出约12\%的题目\scite{nlp-llm-mtp2024}。推断时可以只保留第一个头，按普通方式生成；也可以用其余头提出候选，再由第一个头验证，这属于第\ref{chap:nlp-model-inference}章的推测解码。

独立头有一个结构性限制：第$k$个头只以$x_{1:t}$为条件，看不到中间的$k-1$个词元，各头给出的是给定同一前缀的边缘分布，而不是未来几个词元的联合分布。设在“读书会改到”之后，数据中以各0.5的概率出现“周五下午”或“周六上午”（教学设定）。理想的第一个头给“周五”与“周六”各0.5的概率，第二个头给“下午”与“上午”各0.5的概率。若把两个边缘分布相乘作为联合分布，“周五上午”会得到0.25的概率，而它在数据中从未出现。因此独立头适合作为辅助训练信号，若要用它们直接同时输出多个词元，就必须另有验证步骤。

DeepSeek-V3采用顺序连接的MTP模块，保留了完整的因果链\scite{deepseek2024v3}。设主模型最后一层在位置$i$的表示为$\bm h_i^{0}\in\mathbb R^{d}$。第$k$个模块（$k=1,\ldots,D$）把前一深度的表示与位置$i+k$真实词元的嵌入拼接，再线性投影回$d$维：
\begin{equation}
 \bm h_i'^{\,k}=\bm M_k\bigl[\operatorname{RMSNorm}(\bm h_i^{k-1});\,
 \operatorname{RMSNorm}(\operatorname{Emb}(x_{i+k}))\bigr],
 \qquad \bm M_k\in\mathbb R^{d\times 2d},
 \label{eq:nlp-llm-mtp-sequential}
\end{equation}
其中$[\cdot;\cdot]$表示拼接成$2d$维向量。投影后的序列$\bm h_{1:T-k}'^{\,k}$经过一个带因果掩码的Transformer块$\operatorname{TRM}_k$得到$\bm h_{1:T-k}^{k}$，再经输出头给出$x_{i+k+1}$的分布：$p^{k}_{i+k+1}=\operatorname{softmax}(\bm U\bm h_i^{k})$。嵌入层$\operatorname{Emb}$与输出矩阵$\bm U$都与主模型共享，每个模块只新增$\bm M_k$和一个Transformer块。第$k$个模块的损失是位置$k+2$到$T$上的平均交叉熵，记为$\mathcal L_{\mathrm{MTP}}^{k}$，总的辅助损失为$\frac{\lambda}{D}\sum_{k=1}^{D}\mathcal L_{\mathrm{MTP}}^{k}$，与主模型的下一词元损失相加。DeepSeek-V3取$D=1$，即额外预测再下一个词元，$\lambda$在训练前段取0.3、后段降为0.1。

与独立头相比，顺序模块的差别在于条件。由于$\operatorname{TRM}_k$是因果的，模块$k$在位置$i$可以读取所有$j\leq i$位置投影后的表示，其中包含$x_{1:i+k}$的嵌入，因此它预测$x_{i+k+1}$时的条件是完整前缀$x_{1:i+k}$。这与下一词元预测的条件相同，区别在于模块$k$只有一层，它必须主要依靠主模型较早位置的表示$\bm h_i^{k-1}$，于是迫使主模型的表示为更远的未来做准备。推断时，MTP模块可以直接丢弃，主模型照常工作；也可以在主模型给出$x_{i+1}$之后，由模块提出$x_{i+2}$的候选，交给主模型验证。

两种多词元预测都只增加训练时的辅助信号，不改变主模型的概率定义。它们的代价是额外的计算与显存，以及权重$\lambda$需要调节：辅助损失过大可能挤占主目标。收益随规模与任务而变，原论文在较小模型上也观察到下降，因此不能认为增加预测头就会普遍改善所有任务。

\subsection{基于去噪重建的目标}
\label{sec:nlp-llm-objective-denoising}

去噪目标先按某种规则损坏输入，再要求模型恢复原文。与自回归目标相比，被损坏位置两侧的内容都可以作为条件，模型因而能学习利用双侧上下文。各种方案的差别可以沿四个维度描述：损坏的单位是词元、跨度还是句子；损坏比例多大；恢复对象是只恢复缺失部分还是整段原文；以及与之配套的是编码器、编码器—解码器还是带前缀的解码器。形式上，设损坏函数$c$把$x$变成$\tilde x=c(x)$，恢复目标为序列$\bm y$，损失是$-\log p_\theta(\bm y\mid\tilde x)$。下面从第\ref{chap:nlp-classical-models}章的遮蔽语言模型出发，依次写出各方案的损坏规则与目标，最后说明UL2怎样把它们组合起来。

\subsubsection{遮蔽语言模型的一般化}
\label{sec:nlp-llm-objective-mlm}

第\ref{sec:nlp-classical-bert}节介绍的BERT使用遮蔽语言模型（masked language model, MLM）：随机选取约15\%的位置，把其中大部分替换为掩码符号、少部分替换为随机词元或保持不变，再用双向编码器预测这些位置的原词元。设被选位置集合为$\mathcal M$，损坏后的序列为$\tilde x$，目标为
\[
 \mathcal L_{\mathrm{MLM}}=-\sum_{i\in\mathcal M}\log p_\theta(x_i\mid\tilde x).
\]
在三元组中，$\mathcal S=\mathcal M$，每个位置的可见集合是全部位置（图~\ref{fig:nlp-llm-objective-masks}(c)），但$\mathcal M$中位置的原内容已不可见。这个式子还隐含一个假设：多个被遮位置在给定$\tilde x$时独立预测。若同时遮住“周五”和“下午”，模型分别预测两者，并不建模它们之间的搭配。

把MLM放进统一框架后，它的两个局限变得清楚。第一，只有$|\mathcal M|/T\approx15\%$的位置计分，同样长的文本提供的监督信号约为下一词元预测的七分之一；ELECTRA改为让判别器判断每个位置的词元是否被替换，使全部位置参与计分，从而提高样本效率，详见第\ref{chap:nlp-classical-models}章。第二，双向编码器在每个位置读取整段输入，它定义的是给定其余位置时的条件分布，而不是序列的从左到右分解，因此不能直接用于生成。为了在保留双侧条件的同时获得生成能力，下面的生成式去噪把恢复交给自回归解码器；第\ref{sec:nlp-llm-objective-diffusion}小节的掩码扩散则把遮蔽比例随机化，使MLM式的恢复本身成为生成过程。

\subsubsection{T5跨度损坏与BART去噪}
\label{sec:nlp-llm-objective-t5-bart}

MLM逐个遮蔽词元，而语言中有意义的单位常是连续的短语，例如“下午两点”；逐词元遮蔽时，模型常能由相邻未遮词元轻易猜出答案。T5的\term{跨度损坏}（span corruption）改为遮蔽连续跨度，并把恢复交给编码器—解码器中的解码器\scite{raffel2020t5}。设序列长度为$T$，损坏率为$r$，平均跨度长度为$\mu$，T5的默认取值为$r=0.15$、$\mu=3$。构造步骤为：

\begin{enumerate}
  \item 计算被损坏的词元数约$rT$，跨度数约$n=rT/\mu$，随机划分各跨度长度，并在序列中随机放置这$n$个互不重叠的跨度；
  \item 在输入中把第$k$个跨度整体替换为一个\term{哨兵词元}（sentinel token）$\langle X_k\rangle$，不同跨度使用不同的哨兵，哨兵标明缺口的位置与编号；
  \item 目标序列依次写出“哨兵—被删跨度”，最后以一个额外的哨兵结束。
\end{enumerate}

以教学序列为例（为便于显示，这里遮蔽的比例高于15\%），选中跨度“周五”和“两点”，得到
\begin{center}
 \begin{tabular}{@{}ll@{}}
  \toprule
  编码器输入 & 读书会 改到 $\langle X_1\rangle$ 下午 $\langle X_2\rangle$ 举行\\
  解码器目标 & $\langle X_1\rangle$ 周五 $\langle X_2\rangle$ 两点 $\langle X_3\rangle$\\
  \bottomrule
 \end{tabular}
\end{center}
编码器对损坏后的输入做双向自注意力，解码器以教师强制自回归地生成目标序列，并在每层经交叉注意力读取编码器输出。损失只在目标序列上计算。

跨度损坏的一个直接好处是目标序列短。取$T=512$、$r=0.15$、$\mu=3$，约77个词元被删除，形成约26个跨度；编码器输入约为$512-77+26=461$个词元，解码器目标约为$77+26+1=104$个词元。与要求解码器重建全部512个词元相比，解码器计算显著减少。Raffel等在固定的计算预算下比较多种目标，发现跨度损坏的下游表现与其他去噪目标相近而训练更快，因此作为T5的默认目标\scite{raffel2020t5}。

BART采用另一种取舍：解码器重建完整原文，损坏方式也更多样\scite{lewis-etal-2020-bart}。它是一个\term{去噪自编码器}（denoising autoencoder），即把损坏文本编码后再解码回原文的模型。BART比较了五种噪声，以教学序列为例：

\begin{itemize}
  \item 词元掩码：与MLM相同，随机词元换为掩码，如“读书会 改到 [M] 下午 两点 举行”；
  \item 词元删除：随机删除词元而不留标记，如“读书会 改到 下午 两点 举行”，模型还要判断哪里缺了内容；
  \item 文本填充：按泊松分布（$\lambda=3$）抽取跨度长度，每个跨度换为\emph{单个}掩码符号，长度为0的跨度相当于插入一个掩码，如“读书会 改到 [M] 举行”，模型须推断缺失的是几个词元；
  \item 句子置换：按句号切分文档并随机打乱句子顺序，例如把“读书会改到周五下午两点举行。地点仍在三楼。”的两句对调；
  \item 文档旋转：随机选一个词元，把文档旋转为从它开始，如“下午 两点 举行 读书会 改到 周五”，模型须找出原文开头。
\end{itemize}

无论哪种噪声，解码器的目标都是完整原文，损失为$-\sum_{t}\log p_\theta(x_t\mid x_{<t},\tilde x)$。BART的消融显示文本填充最为有效，最终大规模训练组合了文本填充与句子置换：遮蔽每篇文档约30\%的词元，并打乱全部句子\scite{lewis-etal-2020-bart}。

T5与BART共享编码器—解码器结构和“损坏—恢复”的思路，区别在于恢复对象。T5只生成缺失部分，目标短、计算省；BART重建全文，解码器在训练中大量练习复制与改写未损坏的部分，这与摘要、改写等输出与输入高度重合的生成任务更接近。两者都能直接生成，但生成时依赖一段给定的输入；若要无条件续写，就需要把续写也写成“输入—输出”格式，这引出了下面的前缀语言模型。

\subsubsection{前缀语言模型}
\label{sec:nlp-llm-objective-prefix-lm}

编码器—解码器用两组参数分别处理条件和输出。能否只用一个解码器，同时让条件部分双向可见、输出部分因果生成？\term{前缀语言模型}（prefix language model, prefix LM）通过修改注意力掩码做到这一点：把序列分为长度为$P$的前缀与其后的目标，令
\[
 M_{i,j}=\begin{cases}
 1,& j\leq P\ \text{或}\ j\leq i,\\
 0,& \text{其他},
 \end{cases}
\]
并只对$t>P$的位置计分，$w_t=\mathbb 1[t>P]$。前缀内部的位置可以互相读取，目标位置读取全部前缀和此前的目标。对教学序列取$P=2$，掩码即图~\ref{fig:nlp-llm-objective-masks}(b)：“读书会”与“改到”互读，“周五”及之后的位置按因果方式预测。UniLM用同一个Transformer在双向、从左到右和序列到序列三种掩码之间切换，前缀语言模型对应其中的序列到序列掩码\scite{dong2019unilm}；T5的比较实验中也包含这种结构\scite{raffel2020t5}。

比较图~\ref{fig:nlp-llm-objective-masks}(b)与(d)可以看到，前缀语言模型与编码器—解码器的可见模式相同。差别在于前缀语言模型的条件与输出共用一套参数，目标位置在每一层读取前缀在同一层的表示，而编码器—解码器的解码器只读取编码器最后一层的输出。与下一词元预测相比，前缀语言模型的条件表示可以利用双侧上下文；代价是前缀位置不计分，同样长的文本提供的监督信号更少。此外，前缀表示依赖整个前缀，若前缀在对话中不断追加内容，已经算好的前缀表示会因新内容而改变，不能像因果模型那样直接复用键值缓存，这一点在第\ref{sec:nlp-llm-objective-architecture}小节还会用到。

\subsubsection{UL2混合去噪器}
\label{sec:nlp-llm-objective-ul2}

上述目标各有所长：短跨度损坏训练出的模型在微调后的理解类任务上表现好，因果或前缀语言模型更适合开放生成与少样本提示。UL2的思路是不在其中挑选一个，而是用同一个模型同时学习多种去噪任务，称为\term{混合去噪器}（mixture-of-denoisers）\scite{nlp-llm-ul2-2022}。

UL2把跨度损坏用两个参数描述：平均跨度长度$\mu$与损坏率$r$，跨度数约为$rT/\mu$。在此基础上定义三类去噪器。R去噪器（regular）即T5式的常规跨度损坏，取较短跨度和较低损坏率，如$(\mu,r)=(3,0.15)$或$(8,0.15)$，侧重利用局部上下文恢复知识。X去噪器（extreme）取长跨度或高损坏率，如$(3,0.5)$、$(8,0.5)$、$(64,0.15)$、$(64,0.5)$，输入中只剩少量内容，要求模型根据有限条件生成较长文本。S去噪器（sequential）即前缀语言模型：在随机位置把文本切为前缀与后段，只给前缀，恢复后段。训练时按预先设定的比例从这些去噪器中抽取，构造各自的输入与目标。

不同去噪器要求的行为不同，为了在下游使用时指定行为，UL2在输入开头加入\term{模式标记}（paradigm token）：R、S、X三类分别对应$[\mathrm{NLU}]$、$[\mathrm{S2S}]$、$[\mathrm{NLG}]$。微调或少样本推断时，在输入前加入相应标记，即可切换到与任务相符的模式。以教学序列为例：R模式下输入为“$[\mathrm{NLU}]$读书会 改到 $\langle X_1\rangle$ 下午 两点 举行”，目标为“$\langle X_1\rangle$周五”；X模式可以只保留“读书会”而遮去其后的长跨度；S模式输入“$[\mathrm{S2S}]$读书会 改到”，目标为“周五 下午 两点 举行”。

UL2说明，目标之间的选择不必是互斥的。它的混合目标既可用于编码器—解码器，也可用于仅解码器结构；原论文的20B模型采用编码器—解码器。代价是需要设定多种去噪器的参数与混合比例，以及推断时正确选择模式标记；混合带来的收益依赖这些设置，是在其实验范围内的经验结论。

表~\ref{tab:nlp-llm-objective-denoising}把本小节的去噪目标按“解决了前一方法的什么问题”排成一条谱系。从MLM出发，每一步改动一个维度：T5改变损坏单位与恢复结构，BART改变恢复对象与噪声种类，前缀语言模型把编码器—解码器的可见模式搬进单一解码器，UL2把不同的$(\mu,r)$与模式组合起来。

\begin{table*}[htbp]
 \small
 \begin{tabularx}{\linewidth}{@{}p{22mm}XXX@{}}
  \toprule
  方法 & 解决的问题 & 核心机制 & 新的代价\\
  \midrule
  MLM & 利用双侧上下文 & 遮蔽约15\%词元并逐个恢复 & 计分稀疏；不能直接生成\\
  T5跨度损坏 & 词元级遮蔽过易；需要生成 & 跨度换为哨兵，解码器只生成缺失跨度 & 需编码器—解码器两套参数\\
  BART & 恢复对象过窄 & 五种噪声，解码器重建全文 & 目标长，解码计算多\\
  前缀语言模型 & 用单一解码器获得双向条件 & 前缀块内双向，其后因果 & 前缀不计分；追加内容需重算前缀\\
  UL2 & 单一目标偏向某类任务 & 混合R、S、X去噪器并加模式标记 & 混合参数需设定；推断须选模式\\
  \bottomrule
 \end{tabularx}
 \caption{去噪重建目标的谱系。各行相对于上一行改变一个设计维度；“新的代价”指引入该机制后需要额外付出或承担的部分。}
 \label{tab:nlp-llm-objective-denoising}
\end{table*}

\subsection{基于掩码扩散的目标}
\label{sec:nlp-llm-objective-diffusion}

MLM在固定的15\%遮蔽比例下训练，模型只学会在大部分内容已知时补全少数位置，无法从空白开始生成。若让遮蔽比例从0到1任意变化，模型就要学会在任意程度的缺失下恢复原文；遮蔽比例为1时就是从全掩码序列开始生成。\term{掩码扩散语言模型}（masked diffusion language model）沿这一思路，把生成看成逐步揭开掩码的过程，并为训练目标给出似然上的依据。扩散模型的一般理论可参考\BookChapterRef{chap:pgm-generative}{卷六的概率生成模型章}；这里只用到以下结论：规定一个逐步加噪的前向过程，再学习其反向过程，训练目标可以写成负对数似然的一个上界。

掩码扩散的前向过程很简单。令噪声水平$t\in[0,1]$，给定原文$x_0=x$，对每个位置独立地以概率$t$替换为掩码符号$\mathrm M$、以概率$1-t$保留原词元，得到$x_t$。$t=0$时序列完整，$t=1$时全部被遮蔽。模型$p_\theta(\cdot\mid x_t)$是一个不带因果掩码的Transformer，它读取部分遮蔽的整条序列，对每个被遮位置输出原词元的分布。MDLM与LLaDA使用的训练目标为\scite{nlp-llm-mdlm2024,nlp-llm-llada2025}
\begin{equation}
 \mathcal L_{\mathrm{MD}}(\theta)=
 \mathbb E_{t\sim\mathrm{U}(0,1)}\,\mathbb E_{x_0,\,x_t}
 \left[-\frac{1}{t}
 \sum_{i:\,x_t^i=\mathrm M}
 \log p_\theta\bigl(x_0^i\mid x_t\bigr)\right].
 \label{eq:nlp-llm-masked-diffusion}
\end{equation}
期望覆盖数据、噪声水平与随机遮蔽。在线性遮蔽日程（保留概率$1-t$）下，该目标是模型负对数似然$-\log p_\theta(x_0)$的一个上界，因此最小化它可以看作近似的最大似然训练；采用其他遮蔽日程时，权重$1/t$要相应改为由日程导出的系数\scite{nlp-llm-mdlm2024}。

权重$1/t$的作用可以这样理解。在噪声水平$t$下，长度为$T$的序列平均有$tT$个位置被遮蔽，求和项的规模约与$t$成正比；乘以$1/t$后，每个噪声水平的贡献大致相当于$T$个词元的损失，低遮蔽比例的样本不会因为计分位置少而被忽略。以教学数值为例，$T=6$、$t=0.5$时平均遮蔽3个位置，若各位置损失都为0.4，这一项为$2\times(3\times0.4)=2.4$；$t=0.1$时平均遮蔽0.6个位置，权重变为10。

与MLM相比，式\eqref{eq:nlp-llm-masked-diffusion}只改了两处：遮蔽比例从固定的0.15变为随机抽取的$t$，各项乘以$1/t$。正是这两处改动使目标成为似然上界，并让同一个网络能处理从轻微缺失到完全空白的各种输入。实现上，一步训练的过程如下。

\Needspace{12\baselineskip}
\begin{algorithm}[掩码扩散语言模型的一步训练]
  \label{alg:nlp-llm-masked-diffusion}
  \begin{algorithmic}[1]
    \Require 小批次文本$\{x^{(b)}\}_{b=1}^{B}$，每条长度$T$；双向Transformer $p_\theta$
    \Ensure 参数梯度
    \For{$b=1,\ldots,B$}
      \State 抽取$t_b\sim\mathrm U(0,1)$
      \State 对每个位置$i$独立抽取$u_i\sim\mathrm U(0,1)$，若$u_i<t_b$则令$x_t^i=\mathrm M$，否则$x_t^i=x^{(b),i}$
      \State 前向计算$p_\theta(\cdot\mid x_t)$，得到每个位置的词表分布
      \State $\ell_b\gets-\frac{1}{t_bT}\sum_{i:\,x_t^i=\mathrm M}\log p_\theta(x^{(b),i}\mid x_t)$
    \EndFor
    \State 对$\frac1B\sum_b\ell_b$求梯度
  \end{algorithmic}
\end{algorithm}

算法~\ref{alg:nlp-llm-masked-diffusion}中的$1/T$只是把损失换算到每词元尺度，不改变优化方向。LLaDA按此目标从零训练了约8B参数的模型；在监督微调阶段，它只遮蔽回答部分，提示始终保持可见\scite{nlp-llm-llada2025}。Dream 7B则从自回归模型的权重初始化，再转为掩码扩散训练，以利用已有自回归模型学到的知识\scite{nlp-llm-dream2025}。对教学序列，受损输入可以同时保留“周五”和“举行”而要求恢复中间时间；当$t$接近1时，模型几乎只能依靠学到的文本分布。

掩码扩散的优点是能利用双侧上下文、可以一次恢复多个位置，并且生成顺序不必从左到右。它的代价也来自这些性质：模型给出的是被遮位置的恢复分布，不是式\eqref{eq:nlp-llm-sequence-probability}中的前缀条件概率，不能把这些输出连乘得到文本概率，只能用式\eqref{eq:nlp-llm-masked-diffusion}给出上界；因为每一步都对整条序列做双向计算，自回归模型的键值缓存不能直接复用；生成长度通常需要预先给定。从全掩码序列逐步揭开的采样算法在第\ref{chap:nlp-model-inference}章展开。损坏过程、恢复目标和生成过程三者必须配套定义，换用其中一项时其余两项也要随之调整。

\subsection{编码器—解码器与仅解码器结构的选择}
\label{sec:nlp-llm-objective-architecture}

目标与结构相互约束。跨度损坏与BART去噪的输入和输出分开，自然使用编码器—解码器；下一词元预测、中间填充与多词元预测只需要因果解码器；前缀语言模型介于两者之间。本小节先写清编码器—解码器的计算，再讨论当前大语言模型多采用仅解码器结构的原因，并把这些原因限定为工程与经验观察。

\subsubsection{交叉注意力与编码器—解码器计算}
\label{sec:nlp-llm-objective-cross-attention}

\term{编码器—解码器}（encoder--decoder）结构由两个Transformer堆叠组成。编码器读取长度为$S$的输入$\tilde x$，经$L_{\mathrm e}$层双向自注意力与前馈网络，得到$\bm H_{\mathrm{enc}}\in\mathbb R^{S\times d}$。解码器处理长度为$T$的目标序列，每层依次包含三个子层：带因果掩码的自注意力、交叉注意力和前馈网络，各子层带残差连接与归一化。\term{交叉注意力}（cross-attention）的查询来自解码器，键和值来自编码器输出。设解码器某层交叉注意力子层的输入为$\bm H_{\mathrm{dec}}\in\mathbb R^{T\times d}$，单头计算为
\begin{equation}
 \begin{gathered}
 \bm Q=\bm H_{\mathrm{dec}}\bm W_Q,\quad
 \bm K=\bm H_{\mathrm{enc}}\bm W_K,\quad
 \bm V=\bm H_{\mathrm{enc}}\bm W_V,\\
 \operatorname{CrossAttn}=\operatorname{softmax}\!\left(\frac{\bm Q\bm K^{\top}}{\sqrt{d_h}}\right)\bm V,
 \end{gathered}
 \label{eq:nlp-llm-cross-attention}
\end{equation}
其中$\bm W_Q,\bm W_K,\bm W_V\in\mathbb R^{d\times d_h}$，$\bm Q\in\mathbb R^{T\times d_h}$，$\bm K,\bm V\in\mathbb R^{S\times d_h}$，注意力矩阵为$T\times S$，输出为$T\times d_h$。交叉注意力不需要因果掩码，因为编码器输出在生成开始前已经完整给出；若输入含填充，则对填充列加掩码。多头情形与自注意力相同，各头输出拼接后乘以输出投影。

这一结构在生成时有一个便利：编码器只运行一次，每个解码层由$\bm H_{\mathrm{enc}}$算出的$\bm K,\bm V$在整个生成过程中保持不变，可以缓存复用；解码器自注意力的键值则像普通因果模型一样逐步追加。

两种结构的参数与计算分配不同。设编码器与解码器各$L$层、宽度相同，编码器—解码器的参数约为$2L$层仅解码器模型的参数（解码层因交叉注意力略多）。处理长$S$的输入和长$T$的输出时，输入词元只经过编码器的$L$层，输出词元只经过解码器的$L$层，每个词元的计算约相当于一个$L$层仅解码器模型处理$S+T$个词元。Raffel等据此指出，编码器—解码器以约两倍的参数获得与$L$层仅解码器相近的计算量\scite{raffel2020t5}。另一面，编码器—解码器把参数固定地分成“理解输入”和“生成输出”两部分，而仅解码器让同一组参数同时承担两者。

\subsubsection{仅解码器结构成为主流的原因}
\label{sec:nlp-llm-objective-decoder-only}

当前多数通用大语言模型采用\term{仅解码器}（decoder-only）结构，即只用一个带因果掩码的Transformer堆叠，以下一词元预测为主要预训练目标。这个选择没有理论上的最优性证明，它来自以下几方面的工程考虑与经验观察。

第一，统一的接口。几乎所有文本任务都可以写成“给定前文，续写后文”：翻译、问答、摘要都可以把输入放在前面、把输出作为续写。一个续写模型因此能用同一个结构和同一套参数处理理解与生成，并支持在提示中给出示例的少样本用法\scite{brown2020}。第二，训练信号密集。所有位置都参与计分，而MLM与跨度损坏只对少数位置计分，同样的文本与计算下，仅解码器得到更多监督。第三，增量计算方便。多轮对话中新内容不断追加，因果模型的旧位置表示不受后文影响，键值缓存可以直接复用；编码器—解码器和前缀语言模型的条件部分是双向的，追加内容后原有表示会改变，需要重新编码。第四，单一结构便于规模化：只有一种层需要优化与并行，训练与推断系统也只需支持一种计算图。

已有比较实验支持了其中一部分观察，同时说明结论依赖设置。Wang等在相同计算预算下比较了因果解码器、非因果解码器与编码器—解码器，以及多种目标：只做无监督预训练时，以下一词元预测训练的因果解码器零样本表现最好；若再加入多任务提示微调，以掩码类目标训练的非因果结构表现更好\scite{wang2022whatlm}。这说明结构与目标的优劣取决于后续使用方式。在翻译、语音识别等输入与输出明确分离、输入远长于输出的任务中，编码器—解码器仍被广泛使用，例如Whisper\scite{radford2023whisper}。

\subsection{代表目标的对照}
\label{sec:nlp-llm-objective-comparison}

表~\ref{tab:nlp-llm-objective-comparison}按统一描述的各项并列本节的目标。阅读时可以沿列比较：从“可见条件”一列看出谁能利用右侧上下文，从“计分比例”一列看出训练信号的密度，从最后一列看出训练所得模型能否直接用于生成。

\begin{table*}[!b]
 \small
 \begin{tabularx}{\linewidth}{@{}p{20mm}p{30mm}p{26mm}p{18mm}p{22mm}X@{}}
  \toprule
  目标 & 可见条件 & 预测对象 & 计分比例 & 配套结构 & 能否直接生成\\
  \midrule
  下一词元预测 & 左侧前缀 & 每个下一词元 & 全部 & 仅解码器 & 能，从左到右\\
  中间填充 & 前缀与后缀（经重排） & 中间片段（及其余各段） & 全部 & 仅解码器 & 能；另能按格式填空\\
  多词元预测 & 左侧前缀 & 下一及更远词元 & 全部，多头 & 仅解码器加辅助头 & 能；辅助头可丢弃或作草稿\\
  遮蔽语言模型 & 未遮蔽的两侧内容 & 被遮词元 & 约15\% & 编码器 & 不能\\
  T5跨度损坏 & 带哨兵的损坏输入 & 被删跨度 & 约15\%（目标端） & 编码器—解码器 & 能，以输入为条件\\
  BART去噪 & 加噪输入 & 完整原文 & 全部（目标端） & 编码器—解码器 & 能，以输入为条件\\
  前缀语言模型 & 双向前缀与左侧目标 & 前缀后的词元 & 前缀之后 & 仅解码器（改掩码） & 能，以前缀为条件\\
  UL2 & 依去噪器而定 & 跨度或后段 & 依$(\mu,r)$而定 & 两种均可 & 能，按模式标记切换\\
  掩码扩散 & 未遮蔽的两侧内容 & 被遮词元 & 随机比例$t$ & 双向Transformer & 能，迭代揭开掩码\\
  \bottomrule
 \end{tabularx}
 \caption{代表预训练目标的对照。“计分比例”指一个样本中参与损失计算的位置占比；编码器—解码器的目标只在解码器端计分。“能否直接生成”指不经额外改造，训练所得模型能否产生新文本，及其条件形式。}
 \label{tab:nlp-llm-objective-comparison}
\end{table*}

为了看到各目标改变的究竟是什么，表~\ref{tab:nlp-llm-objective-notice}固定教学序列“读书会／改到／周五／下午／两点／举行”，逐一写出各目标的输入、注意力掩码与目标序列。其中遮蔽比例为便于显示而设定，不对应原论文的默认值。

\begin{table*}[htbp]
 \small
 \begin{tabularx}{\linewidth}{@{}p{20mm}p{55mm}p{22mm}X@{}}
  \toprule
  目标 & 输入 & 注意力掩码 & 目标序列（计分部分）\\
  \midrule
  下一词元预测 & BOS 读书会 改到 周五 下午 两点 & 因果 & 读书会 改到 周五 下午 两点 举行\\
  中间填充（PSM） & $\langle\mathrm{PRE}\rangle$读书会 改到$\langle\mathrm{SUF}\rangle$下午 两点 举行$\langle\mathrm{MID}\rangle$周五 & 因果 & 整个序列右移一位，含周五$\langle\mathrm{EOT}\rangle$\\
  多词元预测 & BOS 读书会 改到 周五 下午 两点 & 因果 & 第1头：读书会…举行；第2头：改到…举行\\
  遮蔽语言模型 & 读书会 改到 [M] 下午 两点 举行 & 双向 & 位置3：周五\\
  T5跨度损坏 & 读书会 改到 $\langle X_1\rangle$ 下午 $\langle X_2\rangle$ 举行 & 编码器双向，解码器因果 & $\langle X_1\rangle$ 周五 $\langle X_2\rangle$ 两点 $\langle X_3\rangle$\\
  BART文本填充 & 读书会 改到 [M] 举行 & 编码器双向，解码器因果 & 读书会 改到 周五 下午 两点 举行\\
  前缀语言模型 & 读书会 改到 周五 下午 两点 & 前缀$P=2$ & 周五 下午 两点 举行\\
  掩码扩散（$t=0.5$） & 读书会 [M] 周五 [M] [M] 举行 & 双向 & 改到、下午、两点，权重$1/t=2$\\
  \bottomrule
 \end{tabularx}
 \caption{同一条六词元通知在各目标下的训练样本。各行改变的只是样本构造、注意力掩码与计分位置，网络内部的层运算相同。遮蔽位置与比例为教学设定。}
 \label{tab:nlp-llm-objective-notice}
\end{table*}

两张表合起来回答了本节开头的三个挑战。条件与预测对象由三元组明确写出，因而可以判断各目标拟合的是哪个条件分布；掩码与结构配套写在同一行，可以检查有无泄漏以及能否直接生成；计分比例与可见条件并列，显示了信号密度与双侧上下文之间的取舍。当前大语言模型以下一词元预测为主，并按需要混入中间填充或多词元预测；去噪与掩码扩散目标则在需要双侧条件或非顺序生成时提供了另一类选择。目标确定之后，剩下的问题是在有限预算下怎样优化它，这是第\ref{sec:nlp-llm-optimization}节的内容。
